[Japanese] \documentclass[12pt,a4paper,twocolumn]{article}
\usepackage[UTF8]{ctex}
\usepackage{fontspec}
\usepackage{geometry}
\geometry{left=1.8cm,right=1.8cm,top=2.5cm,bottom=2.5cm,columnsep=0.8cm}
\usepackage{amsmath,amssymb,amsfonts}
\usepackage{graphicx}
\usepackage{booktabs}
\usepackage{multirow}
\usepackage{hyperref}
\usepackage{cite}
\usepackage{setspace}
\onehalfspacing

% 字体设置
\setmainfont{Times New Roman}  % 英文正文字体
\setsansfont{Microsoft YaHei}  % 雅黑（可选）
\setmonofont{Consolas}
\setCJKmainfont{SimSun}        % 宋体
\setCJKsansfont{Microsoft YaHei}

% 标题信息
\title{论文标题}
\author{作者姓名}
\date{\today}

\begin{document}

% 标题页（单栏）
\twocolumn[
\maketitle
\begin{abstract}
\textbf{摘要：} 在此处撰写摘要内容。摘要应简洁明了地概括论文的主要内容、研究方法、主要结果和结论。

\textbf{关键词：} 关键词1；关键词2；关键词3
\end{abstract}
]

% 正文（双栏）
\section{引言}
引言部分应介绍研究背景、研究意义、国内外研究现状以及本文的研究内容和组织结构。

\section{相关工作}
相关工作部分应综述与本文研究相关的已有研究成果，分析现有方法的优缺点，指出本文研究的创新点。

\section{研究方法}
研究方法部分应详细描述本文所采用的研究方法、实验设计、数据来源等。

\subsection{方法概述}
在此处描述方法的基本思路和框架。

\subsection{具体实现}
在此处描述方法的具体实现细节。

\section{实验结果与分析}
实验结果与分析部分应展示实验数据、分析结果，并对结果进行讨论。

\subsection{实验设置}
在此处描述实验的环境、参数设置等。

\subsection{结果分析}
在此处展示和分析实验结果。

\section{结论}
结论部分应总结本文的主要贡献，指出研究的局限性，并提出未来的研究方向。

% 参考文献
\begin{thebibliography}{99}
\bibitem{ref1} 参考文献1
\bibitem{ref2} 参考文献2
\end{thebibliography}

% 附录（可选）
\appendix
\section{附录A}
附录内容（如有需要）。

\end{document}
